\subsection{Flux Landscape with Enhanced Symmetry Not on \(SL(2,\mathbb{Z})\) Elliptic Points --- arXiv:2311.12425}

\textbf{Authors:} Keiya Ishiguro, Takafumi Kai, Tatsuo Kobayashi, Hajime Otsuka

\paragraph{主な主張}
複素構造モジュライが1個のトーラスorientifoldで超対称Minkowski真空を数え上げ、非因子化模型では\(SL(2,\mathbb{Z})\)の楕円点とは異なる、対称性が増強された点に真空の高い縮退が現れることを示した。 \cite{Ishiguro:2023wwf}

\paragraph{新規性}
合同部分群に加えてスケーリング外部自己同型を含む双対性を同定し、基本領域と物理的に異なる真空の数え方を修正した。

\paragraph{位置付け}
フラックス真空の統計とモジュラー対称性の関係を、因子化トーラスから非因子化トーラスへ拡張する研究である。

\paragraph{コメント（読書メモ）}
読書時の理解では、全モジュライ固定に必要な大きなフラックスとtadpole条件との緊張が、tadpole conjectureの問題に当たる（一般に矛盾するとの断定は未確認）。また、特定のフラックスを選ぶ理由という自然さの問題を意識する。ランドスケープの有限性に言及するには、さまざまな弦理論の設定ごとに真空を調べ、双対性で同一視される点を重複して数えないことが必要である。特に対称性が増強される点では、その同一視と縮退を慎重に扱いたい。
